\section{$G_\delta$ diagonals}
\label{sec:small-diagonals}

Theorem~\ref{thm:at-equals-q2} establishes
$\mathfrak{ap}\leq\mathfrak{at}=\mathfrak q_2$.
We now show that the same cardinal is obtained by replacing the Hausdorff
condition with the $T_1$ axiom together with a $G_\delta$-diagonal.

\begin{definition}\label{def:q-diagonal-variants}
For $i\in\{1,2,2\frac{1}{2}\}$, let $\mathfrak q_{i\Delta}$ be the least
cardinality of a second-countable $T_i$ non-$Q$-space with a
$G_\delta$-diagonal. Thus $\mathfrak q_{1\Delta}$,
$\mathfrak q_{2\Delta}$ and $\mathfrak q_{2\frac{1}{2}\Delta}$ require,
respectively, the $T_1$, Hausdorff and Urysohn separation axioms.
\end{definition}

Clearly, for each $i\in\{1,2,2\frac{1}{2}\}$, we have $\mathfrak q_i\leq\mathfrak q_{i\Delta}$.
Moreover, we have
\[
  \mathfrak q_{1\Delta}\leq\mathfrak q_{2\Delta}
  \leq\mathfrak q_{2\frac{1}{2}\Delta}\leq\mathfrak q\leq\mathfrak b.
\]

In this section we will show that $\mathfrak q_{1\Delta}=\mathfrak{at}=\mathfrak q_2$ and that $\mathfrak q_{2\frac{1}{2}\Delta}=\mathfrak q_{2\frac{1}{2}}$.

\begin{lemma}\label{lem:q2-leq-q1delta}
Every second-countable space $X$ with a $G_\delta$-diagonal and
$|X|<\mathfrak q_2$ is a $Q$-space. Consequently,
$\mathfrak q_2\leq\mathfrak q_{1\Delta}$.
\end{lemma}

\begin{proof}
Write
$\Delta_X=\bigcap_{n<\omega}W_n$, where each $W_n$ is open in $X^2$
and $W_{n+1}\subseteq W_n$.
Fix a countable base $\mathcal B$ for $X$. For each $n\in \omega$, let
\[
  \mathcal C_n=\{U\in \mathcal B:U\times U\subseteq W_n\}.
\]
Each $\mathcal C_n$ is an open cover of $X$.
Enumerate $\mathcal C_n=\{U_{n,i}:i<\omega\}$, repeating members
if necessary.
For each $x\in X$, choose $f_x(n)$ with $x\in U_{n,f_x(n)}$.
Since $|X|<\mathfrak q_2\leq\mathfrak q\leq\mathfrak b$,
choose $g\in\omega^\omega$ such that for every $x \in X$, $f_x\leq^*g$. For $N<\omega$, put
\begin{equation}\label{eq:diagonal-finite-cover-pieces}
  X_N=\bigcap_{n\geq N}\bigcup_{i\leq g(n)}U_{n,i}.
\end{equation}
Each $X_N$ is $G_\delta$ in $X$, and $X=\bigcup_{N\in \omega} X_N$.

Fix $N\in\omega$ and write
$V_{n,i}=X_N\cap U_{n,i}$ for $n\geq N$ and $i\leq g(n)$.
Let $\tau_N$ be the original subspace topology on $X_N$, and let
$\sigma_N$ be the topology generated by the sets $X_N\setminus V_{n,i}$.
Their finite intersections form a countable base of $\tau_N$-closed sets.

We claim that $\sigma_N$ is Hausdorff. Given distinct $x,y\in X_N$,
choose $n\geq N$ with $(x,y)\notin W_n$.
Since $U_{n,i}\times U_{n,i}\subseteq W_n$, no $V_{n,i}$ contains both points.
The sets
\[
  H_x=\bigcap_{\substack{i\leq g(n)\\x\notin V_{n,i}}}
       (X_N\setminus V_{n,i}),
  \qquad
  H_y=\bigcap_{\substack{i\leq g(n)\\y\notin V_{n,i}}}
       (X_N\setminus V_{n,i})
\]
are $\sigma_N$-open neighborhoods of $x$ and $y$, respectively.
If $z\in H_x\cap H_y$, choose $i\leq g(n)$ with $z\in V_{n,i}$,
using \eqref{eq:diagonal-finite-cover-pieces}.
The definitions of $H_x$ and $H_y$ then imply $x,y\in V_{n,i}$,
a contradiction. Thus $H_x\cap H_y=\emptyset$.

Fix $A\subseteq X$. Since $|X_N|<\mathfrak q_2$, the space
$(X_N,\sigma_N)$ is a $Q$-space, so $A\cap X_N$ is $G_\delta$
in $\sigma_N$. Every $\sigma_N$-open set is a countable union of
$\tau_N$-closed basic sets. Hence $A\cap X_N$ is $F_{\sigma\delta}$
in $\tau_N$, and we may choose an $F_{\sigma\delta}$ set $B_N$ in $X$
such that $A\cap X_N=B_N\cap X_N$. Since the sets $X_N$ cover $X$, we have
\[
  A=\bigcap_{N\in\omega}\bigl((X\setminus X_N)\cup B_N\bigr).
\]
Each $X\setminus X_N$ is $F_\sigma$, so this expression shows directly
that $A$ is $F_{\sigma\delta}$ in $X$. As $|A|<\mathfrak b$,
Lemma~\ref{lem:diagonal-collapse} implies that $A$ is $F_\sigma$.
Since $A$ was arbitrary, $X$ is a $Q$-space.
\end{proof}

\begin{thm}\label{thm:q1delta-equals-at}
$\mathfrak q_{1\Delta}=\mathfrak{at}=\mathfrak q_2$.
\end{thm}

\begin{proof}
By Lemma~\ref{lem:q2-leq-q1delta} and Theorem~\ref{thm:at-equals-q2},
it remains only to show that $\mathfrak q_{1\Delta}\leq\mathfrak{at}$.

Fix $\kappa<\mathfrak q_{1\Delta}$, an almost disjoint family
$\mathcal A$ of infinite subtrees of an $\omega$-tree $T$ with
$|\mathcal A|\leq\kappa$, and a subfamily $\mathcal Y\subseteq\mathcal A$.
Put $\mathcal Z=\mathcal A\setminus\mathcal Y$.

Choose branches by König's lemma and form $H$, $X$, and $\tau^-$ as in
Lemma~\ref{lem:tree-fsigma-separator}. On the same set $X$, consider
the topology $\tau^+$ generated by
\[
  P_t=\{S\in X:t\in S\}\qquad(t\in T).
\]
Their finite intersections form a countable base. If $S, R\in X$ are distinct,
let $t\in S\setminus R$.
It follows that $S\in P_t$ and $R\notin P_t$,so $(X,\tau^+)$ is $T_1$.
Moreover,
\begin{equation}\label{eq:positive-tree-diagonal}
  \Delta_X
  =\bigcap_{n<\omega}\bigcup_{t\in T_n}(P_t\times P_t).
\end{equation}
Indeed, every infinite subtree meets every level, whereas two distinct
members of $X$ have no common nodes at sufficiently high levels.
Thus $(X,\tau^+)$ has a $G_\delta$-diagonal.

Since $|X|\leq\kappa<\mathfrak q_{1\Delta}$, the space $(X,\tau^+)$
is a $Q$-space, so $H$ is $G_\delta$ in $\tau^+$.
Each $P_t$ is $\tau^-$-closed, because $X\setminus P_t=N_{\{t\}}$.
Hence every $\tau^+$-open set is $F_\sigma$ in $\tau^-$, and $H$ is
$F_{\sigma\delta}$ in $\tau^-$.
As $|H|\leq\kappa<\mathfrak q_{1\Delta}\leq\mathfrak b$,
Lemma~\ref{lem:diagonal-collapse} makes $H$ an $F_\sigma$ set in $\tau^-$.
Since also $|\mathcal Z|<\mathfrak b$, Lemma~\ref{lem:tree-fsigma-separator}
gives a weak separator of $\mathcal Y$ from $\mathcal Z$.
Thus $\kappa<\mathfrak{at}$, proving the remaining inequality.
\end{proof}

We next identify the Urysohn variant.

\begin{lemma}\label{lem:second-countable-urysohn-gdelta-diagonal}
Every second-countable Urysohn space has a $G_\delta$-diagonal.
Consequently,
$\mathfrak q_{2\frac{1}{2}\Delta}=\mathfrak q_{2\frac{1}{2}}$.
\end{lemma}

\begin{proof}
Fix a countable base $\mathcal B$ for a second-countable Urysohn
space $X$. For every pair $U,V\in\mathcal B$ with
$\overline U\cap\overline V=\emptyset$, the family
\[
  \mathcal U_{U,V}=\{X\setminus\overline U,\ X\setminus\overline V\}
\]
is an open cover of $X$. These covers form a countable family
$\mathcal V$. Given distinct $x,y\in X$, the Urysohn axiom gives such
a pair with $x\in U$ and $y\in V$. No member of $\mathcal U_{U,V}$
contains both points. Therefore
\[
  \Delta_X=\bigcap_{\mathcal U\in\mathcal V}
    \bigcup_{O\in\mathcal U}(O\times O).
\]
(The case $|X|\leq1$ is immediate.)
The equality of the cardinals follows because the two definitions
therefore use the same class of spaces.
\end{proof}

Together, these results give
\[
  \mathfrak{ap}\leq\mathfrak q_{1\Delta}
  =\mathfrak{at}=\mathfrak q_2
  \leq\mathfrak q_{2\Delta}
  \leq\mathfrak q_{2\frac{1}{2}\Delta}
  =\mathfrak q_{2\frac{1}{2}}\leq\mathfrak q.
\]
The equality $\mathfrak q_{2\Delta}=\mathfrak{at}$ is not established
by these arguments.

Finally, below the corresponding $Q$-space threshold, even the weaker
separation axioms force a $G_\delta$-diagonal.

\begin{lemma}\label{lem:q1-gdelta-diagonal}
Every second-countable $T_i$ space of cardinality $<\mathfrak q_i$ has a
$G_\delta$-diagonal, for
$i\in\{1,2\}$.
\end{lemma}

\begin{proof}
Let $X$ be a second-countable $T_i$ space such that $|X|<\mathfrak q_i$.
By the definition of $\mathfrak q_i$, the space $X$ is a $Q$-space.  Hence
every subset of $X$ is an $F_\sigma$ subset of $X$.

Fix a countable base $\mathcal B$ for $X$.  For each $B\in\mathcal B$, choose
a countable collection $\mathcal F_B$ of closed subsets of $X$ such that
$B=\bigcup\mathcal F_B$, and put
\[
    \mathcal F=\bigcup_{B\in\mathcal B}\mathcal F_B.
\]

Then $\mathcal F$ is a countable closed network for $X$, meaning that for every
open $O\subseteq X$ and every $x\in O$ there is $F\in\mathcal F$ such that
$x\in F\subseteq O$.
For each pair $(F,B)\in\mathcal F\times\mathcal B$ such that $F\subseteq B$,
let
\[
    \mathcal U_{F,B}=\{B,X\setminus F\}.
\]
This is an open cover of $X$.  Also, the countable family of covers
$\mathcal V=\{\mathcal U_{F,B}:F\in\mathcal F,
      \ B\in\mathcal B,\ F\subseteq B\}$
separates points in the following sense: if $x\neq y$, choose
$B\in\mathcal B$ with $x\in B\subseteq X\setminus\{y\}$, and then choose
$F\in\mathcal F_B$ with $x\in F$; for this pair $(F,B)$, no member of
$\mathcal U_{F,B}$ contains both $x$ and $y$.
Thus,
\[
\Delta_X=\bigcap_{\mathcal U\in\mathcal V}\bigcup_{U \in \mathcal U}U\times U.
\]
\end{proof}
